\documentclass[10pt,a4paper]{amsart}
\usepackage[margin=19mm]{geometry}
\usepackage{amsmath,amssymb,mathtools}
\usepackage[scr=rsfs,cal=euler]{mathalfa}
\usepackage{xfrac}
\usepackage[unicode,hidelinks,bookmarks=false]{hyperref}
\newcommand{\ie}{i.e.\ }
\newcommand{\cf}{cf.\ }
\newcommand{\resp}{respectively\ }
\newcommand{\Subsection}{\subsection}
\providecommand{\nobreakdash}{\nobreak\hbox{-}}
\setlength{\emergencystretch}{2em}
\multlinegap0pt
\usepackage{amssymb}
\usepackage[normalem]{ulem}
\newtheorem{lemma}{Lemma}
\newcommand{\rank}{\mathrm{rank}}
\title{Secret Sharing on Superconcentrator}
\author{Yuan Li}
\date{Source: arXiv:2302.04482v2 (CC BY 4.0)}
\begin{document}
\maketitle

\noindent\emph{Source and licence.} Secret Sharing on Superconcentrator. Version arXiv:2302.04482v2 (CC BY 4.0), distributed by arXiv under the Creative Commons Attribution 4.0 International licence, http://creativecommons.org/licenses/by/4.0/, as recorded in the arXiv licence metadata for this exact version and shown on the abstract page https://arxiv.org/abs/2302.04482v2. Full licence notice: \texttt{licenses/arxiv-2302.04482-CC-BY-4.0.txt}.\par\smallskip

\noindent\textit{Editorial scope.} Counted unit: the article's lemma lem:tm1\_paths (source0.tex lines 479--482). The article's complete original proof of it is delivered with it (source0.tex lines 483--554, one \texttt{proof} environment of 72 lines). No mathematics is written, repaired or re-derived: the statement and the proof are the article's own lines, in the article's own notation, and no retained line is substituted or re-flowed.  No further source-local context is needed: every cross-reference of every retained line resolves inside the two delivered spans.  Nothing was omitted from any retained span, so the package carries no omission record.  No retained line contains a citation, so the wrapper carries no bibliography and no bibliography entry.  No retained line contains a figure environment, an image-inclusion command or drawing code, so no image is delivered and the package declares no asset dependency.  Editorial formatting only, in addition: an AMS \texttt{amsart} preamble that restates the article's own declarations and macros used by the retained lines, namely the environments \texttt{lemma} on separate counters, together with the standard packages those lines need.  The retained environments print in this document as Lemma 1 in order of appearance, whereas the article prints the same items with the numbering of its own full document.  The complete native source of the article as distributed in the arXiv e-print for this exact version is bundled unchanged as \texttt{sources/137159/source0.tex}.
\begin{lemma}
\label{lem:tm1_paths}
	In the above model, if the circuit computes $n$ shares of some $(t, n)$-threshold secret sharing scheme, then for any subset $T$ of outputs of size $t-1$, there are $t-1$ vertex-disjoint paths connecting $T$ with ($t-1$ inputs from) $R$. Moreover, the $(t-1) \times (\ell-1)$ submatrix $M_{T, R}$ has rank $t-1$.
\end{lemma}

\begin{proof}
It suffices to prove the ``moreover'' part. Because if there do not exist $t-1$ vertex-disjoint paths, by Menger's theorem, we know there exists a subset of vertices of size at most $t-2$ whose removal will disconnect $T$ and $R$. Thus, after fixing $s$, all outputs in $T$ can be written as a function (not necessarily linear) in $R$, which implies that the number of distinct outputs $y_T$ is at most $|\mathbb{F}|^{t-2}$ (after fixing $s$). This contradicts to 
$\rank M_{T, R} = t-1$.

Suppose for contradiction that there exists some $T \subseteq [n]$ of size $t-1$ such that $\rank M_{T, R} < t-1$.

\textbf{Case 1:} $M_{T, 1} = \vec{0}$. In other words, the first input $s$ has no impact on outputs $T$. Since $\rank{M_{T, R}} < t-1$ and adding one extra column with all zeros does not change its rank, we have $\rank{M_{T}} < t-1$, where $M_T$ denotes the submatrix of $M$ consisting of all rows in $T$. Let us add one output to $T$, denoted by $T'$, such that $|T'| = t$. By linear algebra,
\[
\rank{M_{T'}} \le \rank{M_{T}} + 1 < t,
\]
which is not of full rank. (Note that $M_{T'}$ is a $t \times \ell$ matrix.) Thus, there exists a subset of $T'$, denoted by $T''$, such that $|T''| < t$ and $\rank{M_{T''}} = \rank{M_{T'}}$, which implies that participants in $T''$ would be able to recover the secret, assuming participants in $T'$ can recover the secret. This contradicts the definition of $(t, n)$-threshold secret sharing scheme.

\textbf{Case 2:} $M_{T, 1} \not= \vec{0}$. Note that
\[
\begin{pmatrix}
y_1 \\
y_2 \\
\vdots \\
y_n
\end{pmatrix}
=
M
\begin{pmatrix}
s \\
r_1 \\
\vdots \\
r_{\ell-1}
\end{pmatrix},
\]
where $M$ is an $n \times \ell$ matrix computed by the arithmetic circuit. Denote by $Y_T$ the vector indexed by rows in $T$. So we have
\begin{eqnarray*}
y_T
& = &
M_{T}
\begin{pmatrix}
s \\
r_1 \\
\vdots \\
r_{\ell-1}
\end{pmatrix} 
 =  M_{T, R} \begin{pmatrix}
r_1 \\
\vdots \\
r_{\ell-1}
\end{pmatrix}
+
sM_{T, 1}
\end{eqnarray*}
We assume for contradiction that $\rank M_{T, R} < t-1$.
%For convenience, denote by $R = \begin{pmatrix}
%r_1 \\
%\vdots \\
%r_{t-1}
%\end{pmatrix}$.

\textbf{Case 2.1:} there exists $X \in \mathbb{F}^{\ell-1}$ such that $M_{T, R}X = M_{T, 1}$. In other words, the first column of $M_{T}$ is a linear combination of the remaining $\ell-1$ columns. By linear algebra, we have
\[
\rank M_{T} = \rank M_{T, R}.
\] 
Consider any superset of $T$ of size $t$, denoted by $T'$. By linear algebra, we have
\[
\rank M_{T'} \le \rank M_{T} + 1 < t.
\]
Using the same argument as Case 1, we know there exists a proper subset of $T'$ of size $t-1$ whose collation can recover the secret.

\textbf{Case 2.2:} there does not exist $X \in \mathbb{F}^{\ell-1}$ such that $M_{T, R}X = M_{T, 1}$. Let $d = \rank M_{T, R}$, which is less than $t-1$ by assumption. Observe that
\begin{itemize}
	\item The image of $X \mapsto M_{T, \{2, 3, \ldots, t\}} X$, where $X \in \mathbb{F}^{\ell-1}$, is a linear subspace of $\mathbb{F}^{t-1}$ of dimension $d$, denoted by $L \subseteq \mathbb{F}^{t-1}$.
	\item For distinct $s \in \mathbb{F}$, $sM_{T, 1} + L$ are disjoint. (Otherwise, $M_{T, 1}$ would be written as a linear combination of the columns in $M_{T, R}$.)
\end{itemize}
So we conclude the collation of $T$ can always recover $s$. Contradiction!
\end{proof}

\end{document}
